\documentclass[letterpaper,12pt]{article}

\usepackage[margin=1in]{geometry}
\setlength\parindent{0pt}
\setlength{\parskip}{0.8em}

\usepackage{titlesec}
\titlespacing*{\section}{0pt}{0.5\baselineskip}{0.2\baselineskip}
\titlespacing*{\subsection}{0pt}{0.4\baselineskip}{0.1\baselineskip}

\usepackage{enumitem}
\setlist{nosep, topsep=2pt}

\usepackage{graphicx}
\graphicspath{{figures/stability/}}

\usepackage{subcaption}
\usepackage{booktabs}
\usepackage{amsmath}
\usepackage{amssymb}
\usepackage[hidelinks]{hyperref}

\title{Stability Analysis --- Lab 3\\
       \large Connecting Language to fMRI, Stat 214, Spring 2026\vspace{-1em}}
\date{}

\begin{document}
\maketitle


\section{Stability Analysis}

In the previous studies we have shown the variation of test results based on different choices of embedding
method, we will then assess the stability of out method based on whether they are robust to the choice
of subject.
In other words, does the ridge regression encoding model capture a real,
generalisable language-to-brain mapping that replicates across individuals,
or is it overfitting to idiosyncratic noise in one subject's data?
This is the central stability question, examined at four complementary levels
of granularity in this part
We apply this test to the \textbf{best-performing model} --- LoRA BERT
mid4$\star$ (layers 5--8 average, $d = 768$, delay-stacked to 3{,}072
features) --- on both subjects \textit{s2} (94{,}251 voxels) and \textit{s3}
(95{,}556 voxels).
Since a conservative test on a
weaker model would understate the replication that is actually achievable, we choose to only perform stability analysis on the best-performing model.



\subsection{Motivation}

If the encoding model captures a real, generalizable language-to-brain
mapping, the pattern of results should replicate across individuals even
though the two subjects differ physiologically and in brain coverage.
Direct voxel-index comparison is invalid: s2 has 94{,}251 voxels and s3
has 95{,}556, each in its own native space with no spatial registration.
We therefore assess replication at four complementary levels that do not
require a common voxel coordinate system, using the saved weight matrix
$W \in \mathbb{R}^{3072 \times n_\text{voxels}}$ from the LoRA mid4
ridge model.

\begin{enumerate}[label=(\arabic*)]
  \item \textbf{Story-level mean CC.}  For each subject, apply the fitted
        weights to the z-scored test embeddings, then compute the mean
        per-voxel Pearson~$r$ restricted to each of the 20 test stories.
        The resulting 20-dimensional profile is subject-count agnostic
        and directly comparable across subjects.
  \item \textbf{Per-story sorted bar chart.}  Visualise both profiles
        simultaneously, sorted by average CC, to identify which stories
        are universally easy or hard for the model to predict.
  \item \textbf{Voxelwise CC distributions.}  Compare the full
        per-voxel CC histograms for s2 and s3 to determine whether the
        s3 advantage is localised (a few high-CC voxels) or a uniform
        population-level shift.
\end{enumerate}

Additionally, we examine how the \textbf{cross-subject CC ratio}
($\overline{r}_{s3}/\overline{r}_{s2}$) varies across all eight embedding
methods to show that the replication pattern is not specific to LoRA mid4.

\subsection{Results}

\paragraph{Story-level replication.}
The headline summary is:

\begin{itemize}
  \item Story-level Pearson $r$ (s2 vs.\ s3, $n = 20$): $\mathbf{0.745}$
  \item s2 mean CC: $0.0374$, median $0.0330$;\quad
        s3 mean CC: $0.0512$, median $0.0432$
  \item s3 $>$ s2 in \textbf{17 of 20} test stories
  \item Top-1\% CC threshold: s2 $= 0.133$, s3 $= 0.180$
  \item Positive voxels: s2 $= 95.5\%$, s3 $= 97.8\%$
\end{itemize}


Figure~\ref{fig:check2_bars} (panel~B) sorts all 20 stories by average CC
and plots s2 and s3 side by side.  Three patterns stand out.
\textit{High for both subjects}:
\textit{onapproachtopluto} (s2 $= 0.075$, s3 $= 0.097$) and
\textit{fromboyhoodtofatherhood} (s2 $= 0.067$, s3 $= 0.098$) are
structurally rich, event-driven narratives whose contextual semantics the
fine-tuned BERT representations capture robustly.
\textit{marryamanwholoveshismother} (s2 $= 0.045$, s3 $= 0.065$) is
similarly well-predicted by both subjects.
\textit{Relatively lower for both}:
\textit{thatthingonmyarm} (s2 $= 0.019$, s3 $= 0.043$),
\textit{goingthelibertyway} (s2 $= 0.022$, s3 $= 0.058$), and
\textit{theadvancedbeginner} (s2 $= 0.020$, s3 $= 0.044$) rank lowest
for s2; notably these stories still produce meaningful CC under the best
model, whereas GloVe yielded near-zero predictions for the weakest stories.
\textit{Reversal stories} (s2 $>$ s3, marked $\star$):
only three stories reverse the typical ordering ---
\textit{reachingoutbetweenthebars} (s2 $= 0.060$, s3 $= 0.046$),
\textit{birthofanation} (s2 $= 0.036$, s3 $= 0.034$), and
\textit{canplanetearthfeedtenbillionpeoplepart1}
(s2 $= 0.028$, s3 $= 0.024$) --- consistent with individual variation
rather than a systematic pattern.

\paragraph{Voxelwise CC distributions.}
Figure~\ref{fig:check2_voxel} (panel~C) shows the full per-voxel CC
histograms for both subjects under the LoRA mid4 model.
Two features stand out compared to GloVe.
First, the distributions are strongly right-shifted: 95.5\% of s2 voxels
and 97.8\% of s3 voxels have positive CC, versus roughly 60--70\% under
GloVe.  The model predicts the direction of response correctly for almost
every voxel.
Second, the s3 distribution is again shifted uniformly rightward, with
the s3 mean (0.051) exceeding the s2 mean (0.037) across the entire
voxel population --- the advantage is not confined to a small elite of
top voxels.


\begin{center}
\begin{tabular}{lcc}
  \toprule
  Method & s2 mean CC & s3/s2 ratio \\
  \midrule
  BoW                  & 0.0069 & $1.99\times$ \\
  Word2Vec             & 0.0124 & $1.87\times$ \\
  GloVe                & 0.0128 & $1.72\times$ \\
  Vanilla BERT         & 0.0218 & $1.62\times$ \\
  Enhanced BERT        & 0.0349 & $1.41\times$ \\
  LoRA (last4/improved)& 0.0344 & $1.42\times$ \\
  LoRA (mid4$\star$)   & 0.0374 & $1.37\times$ \\
  \bottomrule
\end{tabular}
\end{center}

The shrinking ratio means that better embeddings not only raise absolute
CC but also \emph{close the gap between subjects}: richer semantic
representations recover signal that is present in both brains but
suppressed under weaker models.  This convergence is itself a replication
argument --- if the model were capturing idiosyncratic noise in one
subject, the inter-subject ratio would not change systematically with
embedding quality.



\begin{figure}[H]
  \centering
  \includegraphics[width=\linewidth]{check2_mid4_B_story_bars}
  \caption{\textbf{Check 2 (panel B) --- Per-story encoding accuracy for
           both subjects (LoRA mid4$\star$), sorted by mean CC.}  Dashed
           lines mark each subject's global mean CC.  $\star$ = the three
           stories where s2 outperforms s3.  All stories produce positive
           mean CC under the best model; the ranking is highly consistent
           across subjects.}
  \label{fig:check2_bars}
\end{figure}

\begin{figure}[H]
  \centering
    \includegraphics[width=0.8\linewidth]{check2_mid4_C_distributions}
    \caption{Full per-voxel CC distributions for both subjects (LoRA
             mid4$\star$).  95.5\% of s2 voxels and 97.8\% of s3 voxels
             are positive; s3 (red) is shifted uniformly rightward.
             Dashed lines = subject means; dotted lines = top-1\%
             thresholds.}
    \label{fig:check2_voxel}
\end{figure}


\end{document}